\section*{Bab \ref{ch:g8:powers} --- Pangkat}

\begin{solution}{exo:g8:powers:1}
$2^6 = 64$; \quad $3^3 = 27$; \quad $10^5 = 100\,000$; \quad
$1^{100} = 1$; \quad $0.1^2 = 0.01$; \quad $6^0 = 1$.
\end{solution}

\begin{solution}{exo:g8:powers:2}
$(-3)^2 = 9$; \quad $-3^2 = -9$; \quad $(-1)^{15} = -1$; \quad
$(-5)^3 = -125$; \quad $\left(\frac23\right)^2 = \frac49$.
\end{solution}

\begin{solution}{exo:g8:powers:3}
$7^4 \times 7^5 = 7^9$; \quad
$\dfrac{2^9}{2^3} = 2^6$; \quad
$\left(10^3\right)^4 = 10^{12}$; \quad
$5^6 \times 5^{-2} = 5^4$; \quad
$3^4 \times 7^4 = 21^4$ (\omterm{def:g8:powers:def}{eksponennya} sama: kalikan bilangan pokoknya).
\end{solution}

\begin{solution}{exo:g8:powers:4}
$10^{-2} = 0.01$; \quad
$4 \times 10^3 = 4000$; \quad
$2.5 \times 10^{-4} = 0.00025$; \quad
$10^0 = 1$.
\end{solution}

\begin{solution}{exo:g8:powers:5}
$10^5$; \quad $10^{-1}$; \quad $10^{10}$; \quad $10^{-6}$.
\end{solution}

\begin{solution}{exo:g8:powers:6}
$\dfrac{10^7 \times 10^{-3}}{10^2} = \dfrac{10^4}{10^2} = 10^2 =
100$.

$\dfrac{2^5 \times 2^4}{2^6} = \dfrac{2^9}{2^6} = 2^3 = 8$.

$\left(2^2\right)^3 \times 2^{-4} = 2^6 \times 2^{-4} = 2^2 = 4$.
\end{solution}

\begin{solution}{exo:g8:powers:7}
$2^3 \times 2^4 = 2^{12}$: \emph{salah} --- \omterm{def:g8:powers:def}{eksponennya} dijumlahkan:
$2^7$.

$5^2 + 5^3 = 5^5$: \emph{salah} --- tidak ada aturan untuk \omterm{def:g1:addition:def}{jumlah}:
$25 + 125 = 150$, sedangkan $5^5 = 3125$.

$(3^2)^4 = 3^8$: \emph{benar}.

$10^3 \times 10^3 = 100^3$: \emph{benar} --- keduanya sama dengan $10^6$
(kiri: $10^{3+3}$; kanan: $(10^2)^3$).
\end{solution}

\begin{solution}{exo:g8:powers:8}
Orang baru pada hari $4$: masing-masing dari $3^3 = 27$ orang yang
diberi tahu pada hari 3 memberi tahu tiga orang lain: $3^4 = 81$. Yang
mengetahuinya pada akhir hari 4: $3 + 9 + 27 + 81 = 120$ orang.
\end{solution}

\begin{solution}{exo:g8:powers:9}
\emph{1.} Tebalnya: $2^{10} \times 10^{-4}$ m
$= 1024 \times 10^{-4}$ m $\approx 0.1$ m $= 10$ cm.

\emph{2.} $2^{42} \times 10^{-4} \approx 4.4 \times 10^{12} \times
10^{-4} = 4.4 \times 10^8$ m, lebih besar daripada
$3.8 \times 10^8$ m: sesudah $42$ lipatan (secara teori!), gumpalan
kertasnya akan melewati Bulan.
\end{solution}

\begin{solution}{exo:g8:powers:10}
$2^{10} = 1024$; $10^3 = 1000$; $3^6 = 729$; $5^4 = 625$. Urutannya:
\[
5^4 < 3^6 < 10^3 < 2^{10} .
\]
\end{solution}

\begin{solution}{exo:g8:powers:11}
$2^{100} = \left(2^{10}\right)^{10} = 1024^{10}$ dan
$10^{30} = \left(10^3\right)^{10} = 1000^{10}$. Karena $1024 > 1000$,
mengalikan sepuluh salinan masing-masing menjaga ketaksamaannya:
$2^{100} > 10^{30}$.
\end{solution}

\begin{solution}{pb:g8:powers:1}
\textbf{1.} Butirnya berlipat dua dari petak ke petak bermula dari
$1 = 2^0$: petak $k$ memuat $2^{k-1}$ butir
(\cref{def:g8:powers:def}). Petak $64$ memuat $2^{63}$.

\textbf{2.} $S_1 = 1$, $S_2 = 1 + 2 = 3$, $S_3 = 3 + 4 = 7$,
$S_4 = 7 + 8 = 15$, $S_5 = 15 + 16 = 31$: selalu satu kurang daripada
pangkat $2$ berikutnya ($2$, $4$, $8$, $16$, $32$). Dugaan:
$S_n = 2^n - 1$.

\textbf{3.} Melipatduakan tiap suku $S_n$ menggeser tiap pangkatnya
naik satu ($2 \times 2^{k} = 2^{k+1}$):
\begin{align*}
S_n &= 1 + 2 + 4 + \dots + 2^{n-1}, \\
2 \times S_n &= \phantom{1 + {}} 2 + 4 + \dots + 2^{n-1} + 2^n .
\end{align*}
Dengan mengurangkan baris pertama dari baris kedua, setiap suku dari
$2$ sampai $2^{n-1}$ muncul pada keduanya lalu saling meniadakan:
\[
2 \times S_n - S_n = 2^n - 1,
\qquad\text{yaitu}\qquad
S_n = 2^n - 1 .
\]

\textbf{4.} Seluruhnya $S_{64} = 2^{64} - 1$ butir. Petak keenam
puluh lima akan memuat $2^{64}$ butir: seluruh papannya membawa tepat
satu butir lebih sedikit daripada yang akan dimuat satu petak itu.

\textbf{5.} Separuh pertamanya memuat $S_{32} = 2^{32} - 1$ butir.
Seluruh papannya memuat $2^{64} - 1$, jadi separuh keduanya memuat
\[
\left(2^{64} - 1\right) - \left(2^{32} - 1\right)
= 2^{64} - 2^{32}
= 2^{32} \times \left(2^{32} - 1\right)
\]
(faktorkan $2^{32}$, memakai $2^{32} \times 2^{32} = 2^{64}$,
\cref{thm:g8:powers:rules}): tepat $2^{32}$ kali separuh
pertamanya.

\textbf{6.} Menurut aturan \omterm{def:g8:powers:def}{eksponennya}, $2^{64} = 2^{4 + 60} =
2^4 \times \left(2^{10}\right)^6$. Karena $2^{10} = 1024 > 10^3$,
\[
2^{64} > 16 \times \left(10^3\right)^6 = 16 \times 10^{18}
= 1.6 \times 10^{19} .
\]

\textbf{7.} Lebih dari $1.6 \times 10^{19}$ butir yang masing-masing
$5 \times 10^{-2}$ g:
\[
1.6 \times 10^{19} \times 5 \times 10^{-2}
= 8 \times 10^{17} \text{ g} .
\]
Dengan membaginya dengan $10^6$ g tiap ton: lebih dari
$8 \times 10^{11}$ ton --- delapan ratus miliar ton.

\textbf{8.} $\dfrac{8 \times 10^{11}}{8 \times 10^{8}} = 10^3$:
rajanya menjanjikan sedikitnya \emph{seribu tahun} panen seluruh
dunia pada masa kini. (Legendanya bercerita bahwa para penasihatnya
memberi tahu dia demikian.)

\textbf{9.} $2^{19} = 2^9 \times 2^{10} = 512 \times 1024 =
524\,288 < 10^6$, sedangkan $2^{20} = \left(2^{10}\right)^2 =
1024^2 = 1\,048\,576 > 10^6$. Jadi \omterm{def:g8:powers:def}{eksponen} terkecilnya adalah
$n = 20$: dua puluh kali pelipatduaan melampaui satu juta.

\textbf{10.} Bayaran pada hari $n$ adalah $2^{n-1}$ sen (hari 1:
$2^0 = 1$). Sekarang
\[
2^{26} = 2^6 \times \left(2^{10}\right)^2
= 64 \times 1\,048\,576 = 67\,108\,864 < 10^8,
\]
\[
2^{27} = 2 \times 2^{26} = 134\,217\,728 > 10^8 .
\]
Jadi bayaran hariannya pertama kali melampaui $10^8$ sen ketika
$n - 1 = 27$: yaitu pada hari $28$.

\textbf{11.} $21 = 16 + 4 + 1$: anak timbangan $16$, $4$ dan $1$ g.
$27 = 16 + 8 + 2 + 1$: anak timbangan $16$, $8$, $2$ dan $1$ g.

\textbf{12.} Menurut pertanyaan 3, anak timbangan $1, 2, 4, 8$
bersama-sama beratnya $S_4 = 2^4 - 1 = 15$ g. Jadi tanpa anak
timbangan $16$ g pedagangnya tidak dapat melampaui $15$ g: sasaran
$16$ g atau lebih harus memakainya. Dan sasaran $15$ g atau kurang
tidak boleh memakainya, karena anak timbangan $16$ g sendirian sudah
melampaui sasarannya. Karena itu pemilihan anak timbangan yang
terbesar bersifat \emph{terpaksa}. Yang tersisa adalah sasaran paling
banyak $15$ g yang harus disusun dengan $1, 2, 4, 8$ --- dan
penalaran yang sama berulang: $8$ terpaksa (dipakai kalau sasaran
\omterm{def:g3:division:remainder}{sisanya} $\geq 8$, tidak dipakai kalau bukan, karena
$1 + 2 + 4 = 7$), lalu $4$ (karena $1 + 2 = 3$), lalu $2$, lalu $1$.

\textbf{13.} Dengan mengikuti pilihan yang terpaksa itu, sasaran yang
tersisa sesudah tiap tahapnya paling banyak sama dengan \omterm{def:g1:addition:def}{jumlah} anak
timbangan yang tersisa, jadi prosesnya berakhir dengan \omterm{def:g3:division:remainder}{sisa} $0$:
setiap sasaran dari $1$ sampai $31$ tercapai. Dan karena setiap
pilihan sepanjang jalannya terpaksa, tidak ada pilihan anak timbangan
lain yang mencapai sasaran yang sama: penulisan setiap bilangan dari
$1$ sampai $31$ sebagai \omterm{def:g1:addition:def}{jumlah} pangkat $2$ yang berbeda-beda ada dan
\emph{tunggal}.

\textbf{14.} Keenam anak timbangannya berjumlah $S_6 = 2^6 - 1 = 63$
g, dan penalaran pilihan terpaksa yang sama mencakup setiap sasaran
dari $1$ sampai $63$ g. Untuk $45$: sasarannya $\geq 32$, jadi
pakailah $32$; tersisa $13 < 16$, jadi lewati $16$; pakailah $8$
(tersisa $5$), pakailah $4$ (tersisa $1$), lewati $2$, pakailah $1$:
\[
45 = 32 + 8 + 4 + 1 = 2^5 + 2^3 + 2^2 + 2^0 .
\]

\textbf{15.} Memilih $0$ atau $1$ pada masing-masing dari $64$
petaknya berarti memilih pangkat $2^0, 2^1, \dots, 2^{63}$ yang mana
yang dimasukkan ke dalam sebuah \omterm{def:g1:addition:def}{jumlah} --- persis penimbangan
pedagangnya dengan $64$ anak timbangan. Bilangan terkecil yang dapat
disimpan adalah $0$ (semua petaknya bernilai $0$); yang terbesar
adalah \omterm{def:g1:addition:def}{jumlah} \emph{semua} pangkatnya, yaitu total milik rajanya:
$S_{64} = 2^{64} - 1$ (Bagian I). Menurut penalaran pilihan yang
terpaksa, setiap bilangan bulat di antaranya tercapai tepat sekali:
mesin $64$-bit menyimpan persis bilangan bulat dari $0$ sampai
$2^{64} - 1$.

\end{solution}
